\documentclass{beamer}
\usepackage[utf8]{inputenc}

\usetheme{Madrid}
\usecolortheme{default}
\usepackage{amsmath,amssymb,amsfonts,amsthm}
\usepackage{txfonts}
\usepackage{tkz-euclide}
\usepackage{listings}
\usepackage{adjustbox}
\usepackage{array}
\usepackage{tabularx}
\usepackage{gvv}
\usepackage{lmodern}
\usepackage{circuitikz}
\usepackage{tikz}
\usepackage[utf8]{inputenc}

\usetheme{Madrid}
\usecolortheme{default}
\usepackage{amsmath,amssymb,amsfonts,amsthm}
\usepackage{txfonts}
\usepackage{tkz-euclide}
\usepackage{listings}
\usepackage{adjustbox}
\usepackage{array}
\usepackage{tabularx}
\usepackage{gvv}
\usepackage{lmodern}
\usepackage{circuitikz}
\usepackage{tikz}
\usepackage{graphicx}

\setbeamertemplate{page number in head/foot}[totalframenumber]

\usepackage{tcolorbox}
\tcbuselibrary{minted,breakable,xparse,skins}



\definecolor{bg}{gray}{0.95}
\DeclareTCBListing{mintedbox}{O{}m!O{}}{%
  breakable=true,
  listing engine=minted,
  listing only,
  minted language=#2,
  minted style=default,
  minted options={%
    linenos,
    gobble=0,
    breaklines=true,
    breakafter=,,
    fontsize=\small,
    numbersep=8pt,
    #1},
  boxsep=0pt,
  left skip=0pt,
  right skip=0pt,
  left=25pt,
  right=0pt,
  top=3pt,
  bottom=3pt,
  arc=5pt,
  leftrule=0pt,
  rightrule=0pt,
  bottomrule=2pt,

  colback=bg,
  colframe=orange!70,
  enhanced,
  overlay={%
    \begin{tcbclipinterior}
    \fill[orange!20!white] (frame.south west) rectangle ([xshift=20pt]frame.north west);
    \end{tcbclipinterior}},
  #3,
}
\lstset{
    language=C,
    basicstyle=\ttfamily\small,
    keywordstyle=\color{blue},
    stringstyle=\color{orange},
    commentstyle=\color{green!60!black},
    numbers=left,
    numberstyle=\tiny\color{gray},
    breaklines=true,
    showstringspaces=false,
}
%------------------------------------------------------------
%This block of code defines the information to appear in the
%Title page
\title %optional
{2.10.43}
\date{August  2025}
%\subtitle{A short story}

\author % (optional)
{Bhoomika L - EE25BTECH11014}

\begin{document}


\frame{\titlepage}
\begin{frame}{Question:}
 Let $\vec{a} = \vec{i - k}, \vec{b} = x\vec{i} + \vec{j} + (1 - x) \vec{k}$ and $\vec{C} = y\vec{i} + x\vec{j} + (1 + x - y)\vec{k}.$Then $\sbrak{\vec{a}\  \vec{b}\ \vec{c}}$ depends on
\begin{multicols}{}
\begin{enumerate}
\item only x
\item only y
\item Neither x nor y
\item both x and y
\end{enumerate}
\end{multicols}
\end{frame}

\begin{frame}{Row Echelon form}
\begin{align}
\sbrak{\vec{a}\ \vec{b}\ \vec{c}}=\left|\begin{matrix}1&0&-1\\x&1&1-x\\y&x&1+x-y\end{matrix}\right|
 \end{align}
 
\begin{align}
      R_{2} \longrightarrow R_{2}-xR_{1}  \implies \myvec{1 & 0 &-1\\0 & 1 & 1\\y & x & 1+x-y}\\
    R_{3} \longrightarrow R_{3}-yR_{1}\implies \myvec{1 & 0 &-1\\0 & 1 & 1\\0 & x & 1+x}\\
      R_{3} \longrightarrow R_{3}-xR_{2} \implies \myvec{1 & 0 &-1\\0 & 1 & 1\\0 & 0 & 1}
\end{align}
\end{frame}

\begin{frame}{Solution}
\begin{align}
    \sbrak{\Vec{a} \ \Vec{b} \ \Vec{c}}= \brak{1}\left|\begin{matrix}1&1\\0&1\end{matrix}\right|-\brak{0}\left|\begin{matrix}0&1\\0&1\end{matrix}\right|+\brak{-1}\left|\begin{matrix}0&0\\0&1\end{matrix}\right|=1\\
\implies \sbrak{\Vec{a} \ \Vec{b} \ \Vec{c}}= 1
\end{align}
\begin{center}
    $\therefore$ \sbrak{\Vec{a} \ \Vec{b} \ \Vec{c}} is neither dependent on x and y.
\end{center}    
\end{frame}

\begin{frame}[fragile]
    \frametitle{C Code }
    \begin{lstlisting}
    // code5.c
double dot(double a[3], double b[3]) {
    return a[0]*b[0] + a[1]*b[1] + a[2]*b[2];
}

void cross(double u[3], double v[3], double res[3]) {
    res[0] = u[1]*v[2] - u[2]*v[1];
    res[1] = u[2]*v[0] - u[0]*v[2];
    res[2] = u[0]*v[1] - u[1]*v[0];
}
\end{lstlisting}
\end{frame}

\begin{frame}[fragile]
    \frametitle{Python Code }
    \begin{lstlisting}
 import ctypes

# Load the shared library (adjust filename & path if needed)
lib = ctypes.CDLL('./code5.so')  # use 'vectorops.dll' on Windows

# Declare argument and return types for dot
lib.dot.argtypes = [ctypes.POINTER(ctypes.c_double), ctypes.POINTER(ctypes.c_double)]
lib.dot.restype = ctypes.c_double

# Declare argument types for cross
lib.cross.argtypes = [ctypes.POINTER(ctypes.c_double), ctypes.POINTER(ctypes.c_double), ctypes.POINTER(ctypes.c_double)]
lib.cross.restype = None
\end{lstlisting}
\end{frame}

\begin{frame}[fragile]
    \frametitle{Python Code }
    \begin{lstlisting}
def dot(a, b):
    a_c = (ctypes.c_double * 3)(*a)
    b_c = (ctypes.c_double * 3)(*b)
    return lib.dot(a_c, b_c)

    def cross(u, v):
    u_c = (ctypes.c_double * 3)(*u)
    v_c = (ctypes.c_double * 3)(*v)
    res_c = (ctypes.c_double * 3)()
    lib.cross(u_c, v_c, res_c)
    return list(res_c)

# Vector a = i - k
a = [1, 0, -1]
\end{lstlisting}
\end{frame}

\begin{frame}[fragile]
    \frametitle{Python Code }
    \begin{lstlisting}
# Test values for x and y
xs = [0, 1, 2]
ys = [0, 1, -1]

print("x    y    [a b c]")
for x in xs:
    for y in ys:
        b = [x, 1, 1 - x]
        c = [y, x, 1 + x - y]
        bc = cross(b, c)
        scalar_triple = dot(a, bc)
        print(f"{x:<4} {y:<4} {scalar_triple:.2f}"
\end{lstlisting}
\end{frame}
 
\end{document}
